\documentclass[11pt,article,oneside]{memoir}
\usepackage[utf8]{inputenc}
\usepackage[T1]{fontenc}
\usepackage{url}
\usepackage[hidelinks,colorlinks=true,linkcolor=blue]{hyperref}
\usepackage{amsmath,amssymb}
\usepackage{graphicx}
\usepackage{xcolor}
\usepackage{tabularx}
\usepackage{listings}
\usepackage{geometry}
\usepackage{titlesec}
\usepackage{libertine}

\geometry{margin=1in}

\titleformat{\section}{\Large\bfseries\color{darkgray}}{\thesection}{1em}{}
\titleformat{\subsection}{\large\bfseries\color{blue!80!black}}{\thesubsection}{1em}{}

\lstset{
  basicstyle=\small\ttfamily,
  breaklines=true,
  frame=lines,
  backgroundcolor=\color{gray!10}
}

\begin{document}

\begin{center}
    \Huge \textbf{Api Reference} \\
    \vspace{0.5em}
    \large \textit{A CEREBRUM Research Publication (v1.2.4)} \\
    \vspace{1em}
    \large \textbf{Bryan Alexander Buchorn} \\
    Independent Researcher \\
    \vspace{2em}
\end{center}

\begin{abstract}
\textbf{Overview:} This guide provides a conceptual entry point into the CEREBRUM framework. It is designed for practitioners and researchers seeking to understand the core reasoning signals of Framework Implementation Guide.
\end{abstract}

\section{CEREBRUM REST API Reference}

\textbf{Base URL}: \texttt{http://localhost:8200}
\textbf{API Version}: v1.1.0
\textbf{Authen\-tication}: JWT Bearer token (all endpoints except \texttt{/health})

---

\subsection{Authen\-tication}

All endpoints (except \texttt{GET /health}) require a JWT Bearer token:

\begin{lstlisting}
Authorization: Bearer <jwt_token>
\end{lstlisting}

Tokens are generated using HMAC-SHA256 with the \texttt{CEREBRUM\_JWT\_SECRET} environment variable. Token lifetime defaults to 3600 seconds; configure via \texttt{CEREBRUM\_JWT\_TTL}.

\textbf{Error response (401 Unauth\-orized):}
\begin{lstlisting}
{"detail": "Missing or invalid authorization token"}
\end{lstlisting}

---

\subsection{Endpoints}

\subsubsection{Core Reasoning}

\paragraph{\texttt{POST /query}}
Execute a multi-hop reasoning query against the loaded graph.

\textbf{Request body:}
\begin{lstlisting}
{
    "entity": "Marie Curie",
    "max_hops": 3,
    "beam_width": 5,
    "top_k": 10,
    "probabilistic": false,
    "warm_start_strength": 0.0,
    "community_snapshot": true
}
\end{lstlisting}

\begin{center}
{\footnotesize
\begin{tabularx}{\columnwidth}{|>{\raggedright\arraybackslash}X|>{\raggedright\arraybackslash}X|>{\raggedright\arraybackslash}X|>{\raggedright\arraybackslash}X|}
\hline
Field & Type & Default & Description \\ \hline
\texttt{entity} & string & required & Starting entity for traversal \\ \hline
\texttt{max\_hops} & int & 3 & Maximum reasoning depth \\ \hline
\texttt{beam\_width} & int & 5 & Number of candidates per hop \\ \hline
\texttt{top\_k} & int & 10 & Number of answer paths to return \\ \hline
\texttt{probabilistic} & bool & false & Enable Bayesian beam search \\ \hline
\texttt{warm\_start\_strength} & float & 0.0 & First-hop Beta prior seeding (Phase 19) \\ \hline
\texttt{community\_snapshot} & bool & true & Enable query snapshot isolation (Phase 20) \\ \hline
\end{tabularx}
}
\end{center}

\textbf{Response (200 OK):}
\begin{lstlisting}
{
    "query_entity": "Marie Curie",
    "answers": [
        {
            "entity": "Radioactivity",
            "score": 0.847,
            "path": ["Marie Curie", "discovered", "Polonium", "property_of", "Radioactivity"],
            "path_confidence": 0.823,
            "confidence_interval": [0.78, 0.87],
            "communities": [2, 2, 5],
            "hmac": "sha256:a3f4b2..."
        }
    ],
    "traversal_ms": 6.3,
    "hops_explored": 847,
    "snapshot_id": "snap_1743000000"
}
\end{lstlisting}

\textbf{curl example:}
\begin{lstlisting}
curl -X POST http://localhost:8200/query \
  -H "Authorization: Bearer $TOKEN" \
  -H "Content-Type: application/json" \
  -d '{"entity": "Marie Curie", "max_hops": 3}'
\end{lstlisting}

---

\paragraph{\texttt{GET /query}}
Simple GET-style query for integ\-rations that cannot send POST bodies.

\textbf{Query parameters:} \texttt{entity} (required), \texttt{max\_hops} (default: 3), \texttt{top\_k} (default: 10)

\textbf{Response:} Same schema as \texttt{POST /query}.

---

\subsubsection{Graph Information}

\paragraph{\texttt{GET /communities}}
Return the current community partition of the loaded graph.

\textbf{Response (200 OK):}
\begin{lstlisting}
{
    "algorithm": "DSCF",
    "modularity_q": 0.432,
    "num_communities": 7,
    "communities": {
        "0": ["Marie Curie", "Pierre Curie", "Polonium", "Radium"],
        "1": ["Einstein", "Special Relativity", "General Relativity"],
        "2": ["Newton", "Gravity", "Calculus"]
    },
    "computed_at": 1743000000
}
\end{lstlisting}

---

\paragraph{\texttt{GET /graph/stats}}
Return structural statistics for the loaded graph.

\textbf{Response (200 OK):}
\begin{lstlisting}
{
    "num_nodes": 21,
    "num_edges": 30,
    "avg_degree": 2.86,
    "density": 0.143,
    "num_communities": 7,
    "modularity_q": 0.432,
    "avg_clustering_coefficient": 0.41
}
\end{lstlisting}

---

\paragraph{\texttt{GET /graph/entity/{entity\_id}}}
Return all edges and community membership for a specific entity.

\textbf{Path parameter:} \texttt{entity\_id} — URL-encoded entity string

\textbf{Response (200 OK):}
\begin{lstlisting}
{
    "entity": "Marie Curie",
    "community": 0,
    "soft_memberships": {"0": 0.72, "2": 0.18, "5": 0.10},
    "edges_out": [
        {"target": "Polonium", "relation": "discovered", "weight": 0.95},
        {"target": "Radium", "relation": "discovered", "weight": 0.93}
    ],
    "edges_in": [
        {"source": "Nobel_Prize_1903", "relation": "awarded_to", "weight": 0.99}
    ],
    "pagerank": 0.043,
    "betweenness": 0.127,
    "degree": 5
}
\end{lstlisting}

---

\subsubsection{Bridge Twins}

\paragraph{\texttt{GET /bridges}}
Return all active bridge twin nodes.

\textbf{Response (200 OK):}
\begin{lstlisting}
{
    "bridge_count": 3,
    "bridges": [
        {
            "twin_id": "bridge_7",
            "original_id": "Curie_Institute",
            "source_community": 0,
            "destination_community": 3,
            "crossing_count": 47,
            "ltp_weight": 0.84,
            "created_at": 1743000000
        }
    ]
}
\end{lstlisting}

---

\paragraph{\texttt{GET /bridges/{twin\_id}}}
Return details for a specific bridge twin node.

\textbf{Path parameter:} \texttt{twin\_id} — bridge twin identifier

\textbf{Response:} Single bridge object (same schema as items in \texttt{/bridges}).

---

\subsubsection{Streaming}

\paragraph{\texttt{GET /stream/events}}
Server-Sent Event stream of raw graph update events.

\textbf{Headers:}
\begin{lstlisting}
Accept: text/event-stream
Authorization: Bearer <token>
\end{lstlisting}

\textbf{Event format:}
\begin{lstlisting}
event: edge_added
data: {"u": "A", "v": "B", "relation": "CAUSES", "weight": 0.72, "timestamp": 1743000000.0}

event: community_rebalanced
data: {"communities": 7, "q": 0.43, "pruned_bridges": 2}

event: bridge_formed
data: {"twin_id": "bridge_8", "src_community": 0, "dst_community": 3}
\end{lstlisting}

---

\paragraph{\texttt{GET /stream/insights}}
Server-Sent Event stream of mater\-ialized insight events.

\textbf{Headers:} Same as \texttt{/stream/events}.

\textbf{Event format:}
\begin{lstlisting}
event: insight_link
data: {"source": "A", "target": "B", "score": 0.81, "state": "SPECULATIVE"}

event: insight_verified
data: {"source": "A", "target": "B", "fwd_conf": 0.78, "rev_conf": 0.71}

event: insight_refuted
data: {"source": "A", "target": "B", "reason": "bilateral_probe_failed"}
\end{lstlisting}

---

\paragraph{\texttt{POST /stream/push}}
Push a single event into the streaming ingest pipeline.

\textbf{Request body:}
\begin{lstlisting}
{
    "subject": "sensor_42",
    "predicate": "READS",
    "object": "temperature_chamber_1",
    "timestamp": 1743000000.0,
    "weight": 0.85
}
\end{lstlisting}

\textbf{Response (202 Accepted):}
\begin{lstlisting}
{"queued": true, "queue_depth": 14}
\end{lstlisting}

---

\subsubsection{Admini\-stration}

\paragraph{\texttt{GET /health}}
Health check endpoint. No authe\-ntication required.

\textbf{Response (200 OK):}
\begin{lstlisting}
{
    "status": "healthy",
    "version": "1.1.0",
    "graph_loaded": true,
    "num_nodes": 21,
    "num_edges": 30,
    "uptime_s": 3600
}
\end{lstlisting}

---

\paragraph{\texttt{POST /admin/rebalance}}
Trigger an immediate synchronous DSCF community rebalance. Requires \texttt{scope: admin} in JWT.

\textbf{Response (200 OK):}
\begin{lstlisting}
{
    "rebalanced": true,
    "new_q": 0.448,
    "communities": 7,
    "pruned_bridges": 1,
    "duration_ms": 480
}
\end{lstlisting}

---

\paragraph{\texttt{GET /admin/resource-stats}}
Return current ResourceGovernor metrics.

\textbf{Response (200 OK):}
\begin{lstlisting}
{
    "cpu_percent": 12.4,
    "memory_mb": 384,
    "active_queries": 3,
    "queued_events": 47,
    "rebalance_count": 12,
    "last_rebalance_at": 1743000000
}
\end{lstlisting}

---

\paragraph{\texttt{DELETE /admin/cache}}
Clear the mater\-ialized path cache.

\textbf{Response (200 OK):}
\begin{lstlisting}
{"cleared": true, "entries_removed": 1024}
\end{lstlisting}

---

\subsection{Error Codes}

\begin{center}
{\footnotesize
\begin{tabularx}{\columnwidth}{|>{\raggedright\arraybackslash}X|>{\raggedright\arraybackslash}X|>{\raggedright\arraybackslash}X|}
\hline
HTTP Status & Code & Description \\ \hline
400 & \texttt{INVALID\_ENTITY} & Entity not found in graph \\ \hline
400 & \texttt{INVALID\_PARAMS} & Request parameter validation failed \\ \hline
401 & \texttt{AUTH\_MISSING} & Author\-ization header absent \\ \hline
401 & \texttt{AUTH\_EXPIRED} & JWT token has expired \\ \hline
401 & \texttt{AUTH\_INVALID} & JWT signature invalid \\ \hline
403 & \texttt{SCOPE\_INSUFFICIENT} & Token scope does not permit operation \\ \hline
404 & \texttt{NOT\_FOUND} & Requested resource not found \\ \hline
429 & \texttt{RATE\_LIMITED} & Request rate exceeded; retry after N seconds \\ \hline
500 & \texttt{TRAVERSAL\_FAILED} & Internal traversal error (check logs) \\ \hline
503 & \texttt{GRAPH\_NOT\_LOADED} & No graph loaded; call /admin/load first \\ \hline
\end{tabularx}
}
\end{center}

\textbf{Error response format:}
\begin{lstlisting}
{
    "error": "INVALID_ENTITY",
    "detail": "Entity 'xyz' not found in graph",
    "status": 400
}
\end{lstlisting}

---

\subsection{Rate Limiting}

Default rate limits (confi\-gurable via environment variables):

\begin{center}
\begin{tabularx}{\columnwidth}{|>{\raggedright\arraybackslash}X|>{\raggedright\arraybackslash}X|}
\hline
Endpoint & Default Limit \\ \hline
\texttt{POST /query} & 100 requests/minute \\ \hline
\texttt{GET /query} & 100 requests/minute \\ \hline
\texttt{GET /communities} & 30 requests/minute \\ \hline
\texttt{POST /stream/push} & 1000 requests/minute \\ \hline
\texttt{GET /stream/events} & 10 concurrent connections \\ \hline
\texttt{GET /stream/insights} & 10 concurrent connections \\ \hline
\texttt{POST /admin/*} & 10 requests/minute \\ \hline
\end{tabularx}
\end{center}

Rate limit responses include:
\begin{lstlisting}
Retry-After: 15
X-RateLimit-Limit: 100
X-RateLimit-Remaining: 0
X-RateLimit-Reset: 1743000060
\end{lstlisting}

---

\subsection{SDK Usage}

\begin{lstlisting}
# Python client (using httpx)
import httpx, os

client = httpx.Client(
    base_url="http://localhost:8200",
    headers={"Authorization": f"Bearer {os.environ['CEREBRUM_TOKEN']}"}
)

response = client.post("/query", json={"entity": "Marie Curie", "max_hops": 3})
result = response.json()
print(result["answers"][0]["path"])
\end{lstlisting}

---
\textbf{Copyright © 2026 Bryan Alexander Buchorn. All Rights Reserved.}


\vspace{3em}
\begin{center}
    \small \textit{Project CEREBRUM — March 2026 — Hardened Enterprise Security (v1.2.4)}
\end{center}

\end{document}
